% ============================================================
% EduPlanner – arc42 Architekturdokumentation v1.0
% ============================================================

% --- Encoding & Fonts ---
\usepackage[utf8]{inputenc}
\usepackage[T1]{fontenc}

% babel-german not available; basic latin encoding used

% --- Layout ---
\usepackage[a4paper, top=2.5cm, bottom=2.5cm, left=2.5cm, right=2.5cm, headheight=15pt]{geometry}
\usepackage{parskip}
\setlength{\parindent}{0pt}

% --- Colours ---
\usepackage{xcolor}
\definecolor{darkblue}{HTML}{1F4E79}
\definecolor{midblue}{HTML}{2E75B6}
\definecolor{lightblue}{HTML}{D6E4F0}
\definecolor{tablegrey}{HTML}{F2F2F2}
\definecolor{tableheader}{HTML}{1F4E79}
\definecolor{noteyellow}{HTML}{FFF8DC}
\definecolor{noteborder}{HTML}{F0A500}
\definecolor{notegreen}{HTML}{E8F4E8}
\definecolor{notegreenborder}{HTML}{375623}
\definecolor{notered}{HTML}{FDECEA}
\definecolor{noteredborder}{HTML}{C00000}
\definecolor{codebg}{HTML}{F5F5F5}
\definecolor{codefg}{HTML}{1A1A2E}
\definecolor{grey}{HTML}{595959}

% --- Tables ---
\usepackage{booktabs}
\usepackage{longtable}
\usepackage{array}
\usepackage{tabularx}
\usepackage{colortbl}
\usepackage{multirow}
\renewcommand{\arraystretch}{1.35}
\newcolumntype{L}[1]{>{\raggedright\arraybackslash}p{#1}}
\newcolumntype{C}[1]{>{\centering\arraybackslash}p{#1}}

% --- Graphics & Boxes ---
\usepackage{graphicx}
\graphicspath{{assets/}}
\usepackage{tcolorbox}
\tcbuselibrary{skins, breakable}

% Hinweis-Box (gelb)
\newtcolorbox{hinweis}[1][Hinweis]{
  breakable, enhanced,
  colback=noteyellow, colframe=noteborder,
  fonttitle=\sffamily\bfseries\small, title=#1,
  left=4pt, right=4pt, top=4pt, bottom=4pt,
  fontupper=\small\itshape
}

% Review-Box (grün)
\newtcolorbox{reviewbox}[1][Review v2.1]{
  breakable, enhanced,
  colback=notegreen, colframe=notegreenborder,
  fonttitle=\sffamily\bfseries\small\color{white},
  title=#1,
  coltitle=white,
  attach boxed title to top left={yshift=-2mm, xshift=4mm},
  boxed title style={colback=notegreenborder},
  left=4pt, right=4pt, top=6pt, bottom=4pt,
  fontupper=\small\itshape
}

% Korrektur-Box (rot/orange)
\newtcolorbox{korrektur}[1][Korrektur v1.0]{
  breakable, enhanced,
  colback=notered, colframe=noteredborder,
  fonttitle=\sffamily\bfseries\small,
  title=#1,
  left=4pt, right=4pt, top=4pt, bottom=4pt,
  fontupper=\small\itshape
}

% Lehrhinweis-Box (blau)
\newtcolorbox{lehrhinweis}[1][Lehrhinweis]{
  breakable, enhanced,
  colback=lightblue, colframe=midblue,
  fonttitle=\sffamily\bfseries\small,
  title=#1,
  left=4pt, right=4pt, top=4pt, bottom=4pt,
  fontupper=\small\itshape
}

% --- Code ---
\usepackage{listings}
\lstset{
  backgroundcolor=\color{codebg},
  basicstyle=\ttfamily\footnotesize\color{codefg},
  breaklines=true,
  frame=leftline,
  framerule=2pt,
  rulecolor=\color{midblue},
  xleftmargin=8pt,
  xrightmargin=4pt,
  aboveskip=6pt,
  belowskip=6pt,
  extendedchars=true,
  literate=%
    {ä}{{\"{a}}}1 {ö}{{\"{o}}}1 {ü}{{\"{u}}}1%
    {Ä}{{\"{A}}}1 {Ö}{{\"{O}}}1 {Ü}{{\"{U}}}1%
    {ß}{{\ss{}}}1,
}

% --- Headings ---
\usepackage{titlesec}
\usepackage{needspace}
\setcounter{secnumdepth}{3}  % Nummerierung bis subsubsection

\titleformat{\section}{\sffamily\LARGE\bfseries\color{darkblue}}{\thesection}{1em}{}[\vspace{1pt}\color{midblue}\titlerule]
\titleformat{\subsection}{\sffamily\large\bfseries\color{midblue}}{\thesubsection}{1em}{}
\titleformat{\subsubsection}{\sffamily\normalsize\bfseries\color{grey}}{\thesubsubsection}{1em}{}
\titlespacing{\section}{0pt}{18pt}{8pt}
\titlespacing{\subsection}{0pt}{14pt}{4pt}
\titlespacing{\subsubsection}{0pt}{10pt}{2pt}
\preto\subsubsection{\needspace{4\baselineskip}}

% --- TOC ---
\usepackage{tocloft}
\renewcommand{\cftsecfont}{\sffamily\bfseries\color{darkblue}}
\renewcommand{\cftsubsecfont}{\sffamily\color{midblue}}
\renewcommand{\cftsecpagefont}{\sffamily\bfseries}
\setlength{\cftbeforesecskip}{4pt}

% --- Header/Footer ---
\usepackage{fancyhdr}
\pagestyle{fancy}
\fancyhf{}
\fancyhead[L]{\small\sffamily\color{grey}EduPlanner -- arc42 Architekturdokumentation v1.0 (arc42~v9)}
\fancyhead[R]{\small\sffamily\color{grey}Vertraulich -- Intern (Lehrfassung HFTM)}
\fancyfoot[L]{\small\sffamily\color{grey}SkillBridge AG (fiktiv) \textperiodcentered{} März 2026}
\fancyfoot[R]{\small\sffamily\color{grey}Seite \thepage}
\renewcommand{\headrulewidth}{0.4pt}
\renewcommand{\footrulewidth}{0.4pt}
\renewcommand{\headrule}{\color{midblue}\hrule width\headwidth height\headrulewidth}
\renewcommand{\footrule}{\color{midblue}\hrule width\headwidth height\footrulewidth}

% --- Misc ---
\usepackage{enumitem}
\usepackage{hyperref}
\hypersetup{
  colorlinks=true,
  linkcolor=midblue,
  urlcolor=midblue,
  pdftitle={EduPlanner arc42 v1.0 (arc42 Version 9)},
  pdfauthor={SkillBridge AG (fiktiv)},
}

\usepackage{float}

\interfootnotelinepenalty=10000   % Fussnoten nie umbrechen
% microtype disabled - font expansion requires scalable fonts
\usepackage{pifont}
\usepackage{amssymb}

\usepackage{plantuml}

\usepackage{xltabular}   % = longtable + tabularx

% Checkmark / cross helpers
\newcommand{\yes}{\textcolor{notegreenborder}{\ding{51}}}
\newcommand{\no}{\textcolor{noteredborder}{\ding{55}}}



% Thin rule for visual separation
\newcommand{\sectionrule}{\vspace{4pt}\color{midblue}\hrule\vspace{6pt}}

\usepackage[ngerman]{babel}   % für neue deutsche Rechtschreibung